\begin{abstract}
\noindent
\textbf{Objetivo.} Consolidar la capa de datos distribuida y el pipeline analítico paralelo del
Sistema de Soporte Técnico ISP (equipo ACC), evaluando su rendimiento y tolerancia a fallos. \textbf{Método.} Se fragmentó \texttt{tickets} por rango de \texttt{fecha\_apertura} sobre un
clúster de tres nodos de CockroachDB, verificando completitud, reconstrucción y
disyunción~\cite{ozsu2011principles}; se instrumentó con métricas Prometheus validadas con
\texttt{promtool} y pruebas Testcontainers contra \texttt{SERIALIZABLE} real; y se rediseñó un
pipeline de Spark sobre 600{,}000 incidencias con cinco transformaciones, contrastado mediante
protocolo experimental (Jain~\cite{jain1991art}, Wohlin et al.~\cite{wohlin2012experimentation})
frente a un baseline en pandas.
\textbf{Resultados.} El clúster aplica correctamente aislamiento serializable ante escrituras
concurrentes conflictivas (\texttt{SQLSTATE 40001}), las tres métricas pasan el \emph{linter} de
Prometheus sin advertencias, y el pipeline detecta 380{,}753 incidencias reincidentes (68.3\%)
sobre 557{,}701 filas filtradas. El protocolo estadístico completo (50 corridas) muestra que,
para este volumen de datos, el baseline secuencial en pandas fue significativamente más rápido que
Spark incluso con ocho núcleos ($p \approx 5\times10^{-10}$), aunque Spark sí exhibe un
\emph{speedup} interno real al escalar de uno a ocho núcleos ($p$ de Amdahl $\approx 0.53$). El
clúster tolera la caída de un nodo sin pérdida de datos ni de disponibilidad de escritura (0
errores en 691 filas, pico de latencia P95 de 1.59\,s durante la reelección de líder), y pierde
disponibilidad de escritura durante 62\,s al caer dos de tres nodos, recuperándose de inmediato al
reincorporar un nodo. \textbf{Conclusión.} La fragmentación y la
instrumentación aplicadas son correctas y verificables; el protocolo revela además que el
paralelismo distribuido no es gratuito y solo se justifica a partir de cierta escala de datos.
\end{abstract}

\begin{center}\textbf{Abstract}\end{center}
\begin{abstract}
\noindent
\textbf{Objective.} To consolidate the distributed data layer and the parallel analytics pipeline
of the ISP Technical Support Ticket Management System (team ACC), evaluating their performance
and fault tolerance. \textbf{Method.}
The \texttt{tickets} table was fragmented by range on \texttt{fecha\_apertura} over a three-node
CockroachDB cluster, verifying completeness, reconstruction and
disjointness~\cite{ozsu2011principles}; the service was instrumented with Prometheus metrics
validated with \texttt{promtool} and Testcontainers tests against real \texttt{SERIALIZABLE}
isolation; and a Spark pipeline was redesigned over 600,000 incidents with five transformations,
contrasted via an experimental protocol against a pandas baseline. \textbf{Results.} The cluster correctly enforces
serializable isolation under conflicting concurrent writes (\texttt{SQLSTATE 40001}), all three
required metrics pass the Prometheus linter without warnings, and the pipeline detects 380,753
recurring incidents (68.3\%) out of 557,701 filtered rows. The full statistical protocol (50
trials) shows that, at this data volume, the sequential pandas baseline was significantly faster
than Spark even at eight cores ($p \approx 5\times10^{-10}$), although Spark does show a real
internal speedup when scaling from one to eight cores (Amdahl $p \approx 0.53$). The cluster
tolerates a single-node failure without data loss or write-availability loss (0 errors across 691
rows, peak P95 latency of 1.59\,s during leader re-election), and loses write availability for
62\,s when two of three nodes go down, recovering immediately once a node rejoins. \textbf{Conclusion.} The fragmentation and instrumentation
applied are correct and empirically verifiable; the protocol further reveals that distributed
parallelism is not free and only pays off past a certain data scale.
\end{abstract}
